In other words, @@M0@@ is the @@M1@@-@@M2@@-algebra generated by the generators @@M3@@ subject to the relations (cf. V 4.6):
\begin{enumerate}[label=\textup{\roman*)}]
\item @@M4@@ for @@M5@@;
\item @@M6@@.
\end{enumerate}

We begin with the case where all the @@M7@@ are equal to @@M8@@; the theorem then reads:

\sgaStableTarget{sga6.e06.cor.4-7}\textbf{Corollary 4.7.} \uline{Let} @@M9@@ \uline{be a quasi}-\uline{compact scheme}, @@M10@@ \uline{an} @@M11@@\uline{-Module locally free of rank} @@M12@@, @@M13@@ \uline{the flag bundle of} @@M14@@, and @@M15@@ \uline{the classes of the invertible sheaves} @@M16@@ (\uline{notation of} 4.5). \uline{Then} @@M17@@ \uline{is the} @@M18@@-\uline{algebra generated by the} @@M19@@, \uline{subject to the relations}:
@@M20@@

We proceed by induction on @@M21@@. For @@M22@@, the theorem is trivial. Suppose it has been proved for @@M23@@; we then consider @@M24@@, and on @@M25@@ the exact sequence:
@@M26@@
where @@M27@@ is the structural morphism of @@M28@@. By Lemma 4.2, @@M29@@ is identified with the usual flag bundle @@M30@@ over @@M31@@, the canonical sheaves corresponding under this identification. The diagram:
@@M32@@
shows that @@M33@@ factors through @@M34@@. By the induction hypothesis applied to @@M35@@, and Corollary 4.4, @@M36@@ has as a basis over @@M37@@ the elements of the form:
@@M38@@
moreover, @@M39@@ has as a basis over @@M40@@ the elements @@M41@@; thus @@M42@@ has as a basis over @@M43@@ the elements:
